% ==================================================================
% id:       lim-019
% topic:    גבולות ורציפות של פונקציות
% subtopic: משפטי רציפות (ערך ביניים, ויירשטראס)
% solution: full
% ==================================================================
\textbf{תהא
$f:[a,\infty)\longrightarrow\mathbb{R}$
פונקציה רציפה כך שקיים גבול
$\displaystyle\lim_{x\to\infty}f(x)$.}
\begin{enumerate}
\item
הוכיחו כי
$f$
חסומה ב-
$[a,\infty)$.
\hint{חקו את ההוכחה שסדרה מתכנסת היא חסומה: מהגבול קבלו חסם על קרן, ועל הקטע הסגור שנשאר השתמשו במשפט ויירשטראס.}
\begin{solbox}
רעיון ההוכחה דומה מאוד לזה של הוכחה שכל סדרה מתכנסת היא חסומה, שכבר ראינו בעבר.

נסמן
$L=\displaystyle\lim_{x\to\infty}f(x)$.
מהעובדה ש-
$L=\displaystyle\lim_{x\to\infty}f(x)$,
קיים
$x_0\in[a,\infty)$

כך ש-
$|f(x)-L|<1$
לכל
$x>x_0$,
בפרט
$$L-1<f(x)<L+1$$
לכל
$x\in[x_0+1,\infty)$,
אז בקרן זו חסמנו את
$f$.

נשים לב כי
$f$
רציפה בקטע
$[a,x_0+1]$,
לכן, ממשפט וויירשטרס, קיימים
$m,M$
ממשיים כך ש-
$m\leq f(x)\leq M$
לכל
$x\in[a,x_0+1]$.
לסיכום,
לכל
$x\in[a,\infty)$
מתקיים
$$,\min\{L-1,m\}\leq f(x)\leq\max\{L+1,M\}$$
כלומר,
$f$
חסומה בקרן.
\end{solbox}
\item
האם בהכרח יש ל-
$f$
מקסימום וגם מינימום?
\hint{משפט ויירשטראס דורש קטע סגור, ו-$[a,\infty)$ אינו כזה. אם נראה לכם שהתשובה חיובית נסו להוכיח אותה; אחרת חפשו דוגמה נגדית פשוטה.}
\begin{solbox}
באופן כללי,
$f$
לא חייבת להשיג מקסימום ומינימום שלה. לדוגמא, ניקח
$f(x)=\frac{1}{x}$
בתחום
$[1,\infty)$.
הפונקציה הזאת עומדת בתנאים של השאלה, אך לא קיים לה מינימום בתחום הנדון
(המקסימום קיים).
\end{solbox}
\end{enumerate}
